\section{Auditoría de fuentes: convertir un caso personal en una clase de ingeniería organizacional}

La trayectoria profesional de Joe Sullivan abarca la fiscalía federal de delitos informáticos, la security leadership en empresas de tecnología, el ecosistema de divulgación de vulnerabilidades y proyectos de interés público. Por ello, la clase combina experiencia profesional, el relato de un caso personal, juicios de política pública y recomendaciones de liderazgo. Unos buenos apuntes no pueden mezclar estos contenidos como si fueran un solo tipo de «hecho»: el recuerdo que el ponente tiene de los acontecimientos es una declaración hecha en clase, las afirmaciones de la fiscalía son acusaciones, el jurado y los tribunales se ocupan de la responsabilidad jurídica, y este texto se encarga de extraer mecanismos de ingeniería y de gobernanza reutilizables.

\begin{warningbox}{Límites jurídicos y temporales}
Este texto son apuntes de security engineering y governance; no constituye asesoría jurídica. La clase se grabó en enero de 2025, cuando la apelación de Sullivan aún no se resolvía; el Tribunal de Apelaciones del Noveno Circuito confirmó la condena el 13 de marzo de 2025, el 12 de noviembre de 2025 publicó una opinión modificada y negó la revisión en pleno, y la Suprema Corte de Estados Unidos se negó a admitir la certiorari petition el 29 de junio de 2026. Las expectativas sobre la apelación expresadas en clase solo pueden tomarse como una opinión histórica, no como el estado procesal actual.
\end{warningbox}

\begin{importantbox}{Los tres ciclos de gobernanza de esta clase}
\begin{enumerate}
  \item \textbf{External-report loop}: recibir el reporte externo $\rightarrow$ verificarlo y clasificar su gravedad $\rightarrow$ corregir $\rightarrow$ divulgación coordinada o recompensa;
  \item \textbf{Incident loop}: detectar $\rightarrow$ contener el impacto $\rightarrow$ preservar la evidencia $\rightarrow$ investigar, recuperar y hacer la revisión posterior;
  \item \textbf{Disclosure loop}: hechos técnicos $\rightarrow$ análisis jurídico y de materiality $\rightarrow$ decisión de la alta dirección o del consejo de administración $\rightarrow$ comunicación a reguladores, clientes o al público.
\end{enumerate}
Los tres ciclos pueden detonarse a partir de una misma pista, pero no pueden sustituirse por una misma persona, un mismo ticket ni un mismo mecanismo de incentivos.
\end{importantbox}

\subsection{Un security program no es un «CSO héroe»}

Con frecuencia se evalúa a los líderes de seguridad por su capacidad individual: si son capaces de detectar un ataque, de convencer al CEO o de decidir con rapidez en plena crisis. Sin embargo, lo que en realidad necesita una organización a escala es un \term{security program} (programa de seguridad): capacidades del personal bien definidas, procesos repetibles, plataformas técnicas y una asignación clara de facultades y responsabilidades de gobernanza, de modo que la conducta correcta no dependa de la improvisación de un responsable. De la transición de Sullivan de fiscal a empresas de tecnología como eBay, Facebook y Uber, el tema común más valioso no son los nombres de las empresas, sino que el equipo de seguridad debe evolucionar al mismo ritmo que la complejidad de la organización.

\lecturefigure{01-security-program.png}{Un security program sostenible depende a la vez de people, process, platform y governance.}{Subtítulos locales 00:00--03:20; reelaboración conceptual.}

\begin{knowledgebox}{Lectura de la figura: ninguna de las cuatro capas puede faltar}
Primero observe people, en la esquina superior izquierda: el talento, las guardias, la comunicación y la capacidad de decisión determinan si la organización logra comprender un incidente. Process convierte la intake, el incident command, la disclosure y la revisión posterior en rutas repetibles. Platform aporta identity, telemetry, case management y evidence storage. Governance responde quién tiene authority, quién se somete al challenge y quién firma la divulgación externa. Si solo se construyen herramientas, el resultado es «muchas alertas y pocas decisiones»; si solo se depende de expertos, el resultado es «el experto se va y el sistema pierde la memoria».
\end{knowledgebox}

\teachervoice{La historia profesional de Sullivan reúne el derecho, la ingeniería y la organización en una misma perspectiva. Nota de clase: el valor escaso de un security leader no consiste en conocer cada vulnerabilidad mejor que todos los ingenieros, sino en traducir los hechos técnicos a riesgos jurídicos, de negocio y públicos, y al mismo tiempo asegurar que cada juicio lo emita quien tiene la facultad de responder por él.}

\begin{knowledgebox}{Términos clave: funciones y responsabilidades}
El \term{CSO/CISO} es responsable de las capacidades de seguridad, de la información sobre riesgos y del escalamiento; el \term{general counsel} (consejero jurídico general, GC) es responsable del análisis jurídico de la organización y de la gestión de los abogados; el \term{CEO} responde por las prioridades de la empresa y por las decisiones de negocio importantes; el \term{board/disclosure committee} ejerce la supervisión y atiende las divulgaciones relevantes. Las funciones pueden fusionarse según el tamaño de la empresa, pero debe distinguirse de forma explícita «quién aporta los hechos, quién emite la opinión jurídica y quién asume la decisión de negocio».
\end{knowledgebox}

\subsection{Resumen de la sección}

Esta clase no gira en torno a «quién tiene más razón», sino en torno a un operating system: el reporte externo, la respuesta a incidentes y la gobernanza de la divulgación requieren, cada uno, owner, evidence, decision right y registro. El caso personal solo sirve para comprobar por qué estos mecanismos tienen que existir.
